\chapter{Riwayat Pemilihan Referensi dan Tanggapan Dosen}\label{lam:riwayat}

\section*{A.1 Pemilihan Jurnal Awal}

Jurnal pertama yang dipilih pada Progress~1 adalah \textcite{audebert2020}. Ada empat
alasan:
\begin{enumerate}
  \item \textbf{Sesuai topik mata kuliah.} Paper ini menggabungkan dua modalitas yang
  asalnya sama tetapi sifatnya berbeda: tata letak halaman dan isi tulisan.
  \item \textbf{Datanya terbuka.} Citra Tobacco3482 tersedia di Kaggle, dan teks OCR-nya
  dirilis penulis sebagai QS-OCR, sehingga modalitas teks tidak perlu dibuat sendiri.
  \item \textbf{Arsitekturnya ringan.} MobileNetV2 dan FastText dengan CNN~1D dapat
  dilatih pada GPU laptop 4~GB.
  \item \textbf{Tidak ada kode resmi.} Repositori QS-OCR hanya berisi skrip OCR, sehingga
  seluruh arsitektur harus ditulis ulang dari spesifikasi paper. Ini menuntut pemahaman
  metode, bukan sekadar menjalankan kode orang lain.
\end{enumerate}

\section*{A.2 Tanggapan Dosen atas Progress 1}

\begin{table}[H]
  \centering
  \caption{Tanggapan dosen pengampu dan tindak lanjutnya}
  \label{tab:tanggapan}
  \fontsize{10}{12}\selectfont
  \begin{tabularx}{\textwidth}{|>{\raggedright\arraybackslash}p{5.4cm}|>{\raggedright\arraybackslash}X|}
    \hline
    \textbf{Tanggapan} & \textbf{Tindak lanjut} \\ \hline
    Alur OCR $\rightarrow$ teks dengan Tesseract dan FastText sudah dipahami &
    Dipertahankan sebagai kondisi acuan (QS-OCR $\times$ FastText) \\ \hline
    Eksplorasi metode OCR: EasyOCR, PaddleOCR, Keras-OCR &
    Mesin OCR dijadikan variabel penelitian (RM1). Keras-OCR bersyarat, lihat risiko
    R3 \\ \hline
    Celah penelitian belum ditentukan &
    Dirumuskan pada Bab~\ref{bab:pendahuluan} dan Subbab~\ref{sec:kajian-pustaka}, didukung
    hasil studi pendahuluan \\ \hline
    Kemungkinan eksplorasi pada embedding &
    Embedding dijadikan variabel kedua: FastText dan BERT (RM2) \\ \hline
    Teks yang dideteksi sebagai objek, untuk tahap selanjutnya &
    Kotak posisi teks disimpan sejak langkah OCR; rancangannya pada
    Subbab~\ref{sec:lanjutan} \\ \hline
    Referensi masih tahun 2019 dan berbentuk arsip; minimal tahun 2025 &
    Referensi utama diganti ke \textcite{kashyap2026}. Audebert dikutip dari versi
    terbitnya (Springer, 2020) sebagai baseline yang telah direplikasi \\ \hline
  \end{tabularx}
  \sumber{Catatan tanggapan dosen pengampu.}
\end{table}

Satu hal perlu dicatat soal referensi. Paper Audebert sebenarnya telah terbit di ECML
PKDD 2019 Workshops, Springer CCIS vol.~1167 (2020). Kesalahan ada pada Progress~1, yang
mengutip versi arXiv. Namun tahun terbitnya tetap tidak memenuhi syarat, sehingga posisi
referensi utama tetap diganti.

\section*{A.3 Pencarian Referensi Baru}

Pencarian dilakukan melalui Crossref dan Semantic Scholar dengan empat kriteria:
\begin{enumerate}
  \item terbit di jurnal atau prosiding pada 2025 atau sesudahnya, bukan
  \textit{preprint};
  \item membahas klasifikasi dokumen multimodal teks--citra, perbandingan mesin OCR, atau
  ketahanan embedding terhadap derau OCR;
  \item metadata (DOI, volume, halaman) dapat diverifikasi;
  \item teks lengkapnya dapat dibaca, tidak hanya abstraknya.
\end{enumerate}

Hasilnya dirangkum pada Tabel~\ref{tab:referensi-baru}. Untuk \textcite{larson2025}, yang
dibaca baru abstraknya, karena unduhan teks lengkapnya belum berhasil.

\begin{table}[H]
  \centering
  \caption{Referensi utama dan pendukung yang terbit tahun 2025--2026}
  \label{tab:referensi-baru}
  \fontsize{10}{12}\selectfont
  \begin{tabularx}{\textwidth}{|l|>{\raggedright\arraybackslash}X|}
    \hline
    \textbf{Referensi} & \textbf{Peran dalam penelitian ini} \\ \hline
    \textcite{kashyap2026} & Referensi utama: fusion teks--citra terkini pada
    Tobacco3482 \\ \hline
    \textcite{mironczuk2026} & Pembingkai celah: meta-analisis 139 studi fusion untuk
    klasifikasi dokumen \\ \hline
    \textcite{francis2025} & Perbandingan Tesseract, EasyOCR, PaddleOCR, Keras-OCR, dan
    mesin lain \\ \hline
    \textcite{zhang2025} & Bukti bahwa galat OCR merambat ke tugas lanjutan \\ \hline
    \textcite{michail2025} & Ketahanan embedding terhadap derau OCR \\ \hline
    \textcite{larson2025} & Jalan pintas kode ID pada Tobacco3482 \\ \hline
  \end{tabularx}
  \sumber{Penelusuran kelompok.}
\end{table}

\section*{A.4 Mengapa Kashyap et al. sebagai Referensi Utama}

\textcite{kashyap2026} adalah penerus langsung Audebert: keduanya menggabungkan cabang
citra dan cabang teks OCR, dan keduanya diuji pada Tobacco3482. Kashyap juga
membandingkan hasilnya secara eksplisit dengan Audebert.

Paper ini \textbf{tidak direplikasi}. \textit{Fine-tuning}-nya memakan sekitar 63 jam
pada GPU 16~GB dengan RVL-CDIP, jauh di luar perangkat yang tersedia. Yang diambil dari
paper ini adalah:
\begin{itemize}
  \item embedding BERT;
  \item aturan fusion rata-rata logit;
  \item arah lanjutan yang disebut penulisnya sendiri, yaitu pembobotan modalitas adaptif
  dan kondisi degradasi OCR.
\end{itemize}
Celah yang paling menentukan adalah mesin OCR-nya yang tidak disebutkan. Catatan kritis
lengkapnya ada pada Lampiran~\ref{lam:kashyap}.
